\documentclass[11pt]{article}
\usepackage[a4paper,margin=26mm]{geometry}
\usepackage[T1]{fontenc}
\usepackage{lmodern,amsmath,amssymb,amsthm,mathtools,booktabs,microtype}
\usepackage[dvipsnames]{xcolor}
\usepackage[colorlinks=true,linkcolor=MidnightBlue,urlcolor=MidnightBlue,citecolor=MidnightBlue]{hyperref}
\hypersetup{pdftitle={A uniform bound for complete graph minors},pdfauthor={MohammadTaghi Hajiaghayi and Mahdi JafariRaviz}}
\setlength{\parskip}{4pt}
\setlength{\parindent}{0pt}
\setlength{\emergencystretch}{2em}

\newtheorem{lemma}{Lemma}
\newtheorem*{maintheorem}{Main Theorem}
\newtheorem*{thomasonlemma}{Thomason's lemma}

\DeclareMathOperator{\atanh}{atanh}
\newcommand{\E}{\mathbb E}
\newcommand{\Prb}{\mathbb P}
\newcommand{\floor}[1]{\left\lfloor#1\right\rfloor}
\newcommand{\ceil}[1]{\left\lceil#1\right\rceil}
\newcommand{\dd}{\mathrm d}

\title{\textbf{A uniform bound for complete graph minors}}
\author{MohammadTaghi Hajiaghayi \and Mahdi JafariRaviz}
\date{November 2026}

\begin{document}
\maketitle

\begin{abstract}
Graph Minors XIII states, as a result of Thomason, that for every integer
$p\ge1$ every non-null simple graph $G$ with
$|E(G)|\ge2.68\,p\sqrt{\log_2p}\,|V(G)|$ has a $K_p$ minor. Thomason's paper
gives the constant $2.68$ only for large $p$. We prove that for every integer
$p\ge51$, every non-null simple graph $G$ with
$|E(G)|\ge(16/5)\,p\sqrt{\ln p}\,|V(G)|$ has a $K_p$ minor. Since
$(16/5)\sqrt{\ln2}<2.68$, this gives the statement of Graph Minors XIII for
every $p\ge51$. The proof uses Thomason's density reduction and a random
partition of the vertex set into many small blocks at once. Two certificates,
checked with exact rational arithmetic, cover $51\le p\le10^{44}$, and an
analytic argument covers $\ln p\ge100$.
\end{abstract}

\section{Introduction}\label{sec:intro}

All graphs in this note are finite and simple. For a graph $G$ we write
$v(G)=|V(G)|$ and $e(G)=|E(G)|$. A graph is non-null if it has at least one
vertex, and the density of a non-null graph $G$ is $d(G)=e(G)/v(G)$. We write
$\ln$ for the natural logarithm and $\log_2$ for the logarithm to base $2$.

Graph Minors XIII \cite[(5.5), page~83]{gm13} states the following, as a
result of Thomason \cite{thomason}.
\begin{quote}
\textbf{(5.5)} \emph{Let $p\ge1$ be an integer, and let $G$ be a non-null
simple graph with $|E(G)|\ge2.68\,p(\log_2(p))^{1/2}|V(G)|$. Then $G$ has a
$K_p$-minor.}
\end{quote}
Thomason's paper gives the constant $2.68$ only asymptotically, that is, for
all sufficiently large $p$ \cite[page~264]{thomason}. In this note we prove a
bound of the same form that holds for every $p\ge51$.

\begin{maintheorem}
For every integer $p\ge51$, every non-null graph $G$ with
$d(G)\ge\frac{16}{5}\,p\sqrt{\ln p}$ has a $K_p$ minor.
\end{maintheorem}

The Main Theorem gives (5.5) for every $p\ge51$. Indeed,
$\sqrt{\ln p}=\sqrt{\ln2}\,\sqrt{\log_2p}$, so it suffices that
$(16/5)\sqrt{\ln2}\le2.68=67/25$. We have $\ln2<7/10$, because
$e^{7/10}>1+7/10+(7/10)^2/2+(7/10)^3/6>2$, and therefore
\begin{equation*}
 \left(\frac{16}{5}\right)^2\ln2<\left(\frac{16}{5}\right)^2\frac{7}{10}
 =7.168<7.1824=\left(\frac{67}{25}\right)^2.
\end{equation*}
Numerically, $(16/5)\sqrt{\ln2}=2.6641747557\ldots$, so that
$(16/5)\,p\sqrt{\ln p}=2.6641747557\ldots\,p\sqrt{\log_2p}$.

For $1\le p\le50$, the statement (5.5) is proved in the Lean formalization
only, by a different argument, and this note does not cover these cases. The
note makes no claim about the constant $16/5$ for $p\le50$.

Wood \cite[Theorem~1]{wood} proves an explicit bound of the same form for
every $p\ge2$: every graph of density at least $4.5\,p\sqrt{\ln p}$ has a $K_p$
minor. His proof builds the branch sets one after another, and it adds
vertices to each branch set to make it connected. Here we choose all
candidate branch sets at once, as the blocks of one random partition. We
then discard the disconnected blocks, and enough further blocks to remove
the missing adjacencies, and we bound the expected number of discarded
blocks. No vertex is added to a block, so the proportion of nonneighbors of
a vertex stays the same during the construction.

The proof goes as follows. Section~\ref{sec:reduction} recalls Thomason's
density reduction. It reduces the Main Theorem for $p$ to finding a
$K_{p-2}$ minor in small graphs in which every vertex has few nonneighbors.
Section~\ref{sec:partition} gives the tool for these graphs: a random
partition of the vertices into blocks (Lemmas~\ref{lem:partition}
and~\ref{lem:mixed}), with an explicit lower bound $B$ on the expected
number of blocks that form a complete minor. It remains to show that
$B>p-3$. Section~\ref{sec:finite} checks this for $51\le p\le109$ with a
finite certificate, Section~\ref{sec:medium} for $110\le p\le10^{44}$ with a
certificate that covers whole intervals of $p$ at once, and
Section~\ref{sec:tail} for $\ln p\ge100$ by hand.
Section~\ref{sec:proof} puts the three ranges together, and
Section~\ref{sec:lean} describes how the Lean formalization uses this note.
Appendix~\ref{app:reduction} proves the density reduction.

\section{The density reduction}\label{sec:reduction}

Let $\alpha_0>1$ be the solution of $\alpha_0=1+\ln(2\alpha_0)$. It is
unique, because the function $x-1-\ln(2x)$ has derivative $1-1/x>0$ for
$x>1$, is negative at $x=1$, and tends to infinity. We use the following form
of Lemma~1 of Thomason \cite[page~262]{thomason}.

\begin{thomasonlemma}
Let $p\ge2$ and $c\ge3$ be integers. Suppose that every non-null graph $H$
with
\begin{equation}\label{eq:thomason}
 v(H)\le c\qquad\text{and}\qquad 2\delta(H)-v(H)\ge\floor{c/\alpha_0}-3
\end{equation}
has a $K_{p-2}$ minor. Then every non-null graph of density at least $c$ has
a $K_p$ minor.
\end{thomasonlemma}

Thomason states the lemma with a real parameter $r$, and we take
$r=c/\alpha_0$, so that $\alpha_0r=c$ is an integer. The proof makes two
neighborhood reductions, and the two vertices chosen in them give the two
extra branch sets. Appendix~\ref{app:reduction} gives a complete proof of the
lemma in the form above. In \cite{thomason} the lemma is stated in terms of
an extremal function. Below we apply it only with an integer $c$ strictly less than the
density of the graph, so no question about a boundary case arises if one
uses his formulation instead.

We do not need the value of $\alpha_0$, only the bound $\alpha_0<67/25$. To
prove it, recall that $x-1-\ln(2x)$ is strictly increasing for $x>1$ and
vanishes at $\alpha_0$. At $x=67/25$ it equals $42/25-\ln(134/25)$, which is
positive because
\begin{equation*}
 e^{42/25}>
 \sum_{j=0}^{7}\frac{(42/25)^j}{j!}
 =\frac{134}{25}+\frac{110787181}{30517578125}
 >\frac{134}{25}.
\end{equation*}
So $\alpha_0<67/25$. For an integer $c\ge3$ put
\begin{equation}\label{eq:R}
 R=\floor{\frac{25c}{67}}-1.
\end{equation}

Let $H$ satisfy \eqref{eq:thomason}, let $n=v(H)$, and let
$q=n-1-\delta(H)$, which is the largest number of nonneighbors of a vertex of
$H$. Here and below, a vertex is not counted as a nonneighbor of itself.
Since $c/\alpha_0>25c/67$, the second condition in \eqref{eq:thomason}
implies $2\delta(H)-n\ge R-2$. As $2\delta(H)-n=n-2-2q$, this says that
$n-2q\ge R$. Together with $n\le c$ we get
\begin{equation}\label{eq:parameters}
 R\le n\le c\qquad\text{and}\qquad 0\le q\le\frac{n-R}{2}.
\end{equation}
We call a non-null graph $H$ with $n$ vertices a \emph{reduced graph} (for
$c$) if $R\le n\le c$ and every vertex of $H$ has at most $(n-R)/2$
nonneighbors. If $n=R$ or $n=R+1$, then $q=0$, so a reduced graph of this
order is complete.

Now let $p\ge51$, and put
\begin{equation}\label{eq:D}
 C=\frac{16}{5},\qquad D=Cp\sqrt{\ln p},\qquad c=\ceil{D}-1.
\end{equation}
Then $D-1\le c<D$, and $c\ge3$ since $D>4$. A graph as in the Main Theorem
has density at least $D$, and so at least $c$. By Thomason's lemma and
\eqref{eq:parameters}, the Main Theorem for $p$ follows once every reduced
graph for this $c$ has a $K_{p-2}$ minor. In Section~\ref{sec:finite} the
certificate may instead give any integer $c\ge3$ with $c<D$, and the same
reduction applies. The bound $c\ge D-1$ is not needed there.

\section{A random partition into blocks}\label{sec:partition}

To find a large complete minor in a reduced graph $H$, we split a random set
of vertices into blocks of $k$ vertices each. Every vertex has few
nonneighbors, so a random block is likely to be connected and to dominate
almost all of $H$. Then two blocks are likely to be adjacent. We delete the
blocks that fail, and the remaining blocks form a complete minor.

For a set $S$ of vertices of a graph $H$, let $U(S)=V(H)\setminus N_H[S]$ be
the set of vertices that are neither in $S$ nor adjacent to a vertex of $S$.
We say that $S$ does not dominate these vertices. A $K_m$ model in $H$ is a
family of $m$ pairwise disjoint nonempty sets of vertices, each inducing a
connected subgraph, with an edge between every two of them. A graph has a
$K_m$ minor if and only if it has a $K_m$ model.

\begin{lemma}\label{lem:partition}
Let $H$ be a graph with $n$ vertices in which every vertex has at most $q$
nonneighbors. Let $k$ be an integer with $2\le k<n$, put $N=\floor{n/k}$ and
$\pi=q/n$, and suppose that $N\ge2$. Let $A>0$, and put
\begin{align}
 F&=(2^{k-1}-1)\left(\frac{q}{n-k+1}\right)^{k-1},\qquad
 M=\binom N2\left(\frac{An\pi^k}{n-k}\right)^k,\notag\\
 B&=N(1-F-A^{-1})-M.\label{eq:B}
\end{align}
If $B>m-1$ for an integer $m$, then $H$ has a $K_m$ minor.
\end{lemma}

The three loss terms come from blocks that are disconnected ($NF$), blocks
that do not dominate enough vertices ($N/A$), and pairs of blocks with no
edge between them ($M$).

\begin{proof}
Take a uniformly random ordering of $V(H)$. Let $S_1$ be the set of its first
$k$ vertices, $S_2$ the set of the next $k$ vertices, and so on up to $S_N$.
We call these sets blocks. They are disjoint, and each one on its own is a
uniformly random $k$-subset of $V(H)$.

\emph{Domination.} Fix a block $S_i$ and a vertex $v$. If $v\in U(S_i)$,
then all $k$ vertices of $S_i$ are nonneighbors of $v$, and $v$ has at most
$q$ nonneighbors. So
\begin{equation*}
 \E|U(S_i)|\le n\binom qk\Big/\binom nk\le n\pi^k,
\end{equation*}
where the last step uses $(q-j)/(n-j)\le q/n$ for $0\le j<k$. Call $S_i$
\emph{good} if $|U(S_i)|\le An\pi^k$. If $q>0$, Markov's inequality shows
that $S_i$ is not good with probability at most $1/A$. If $q=0$, then
$U(S_i)$ is empty, so every block is good. In both cases the expected number
of blocks that are not good is at most $N/A$.

\emph{Connectivity.} If $H[S_i]$ is disconnected, then there is a partition
of the $k$ positions of $S_i$ into two nonempty parts with no edge of $H$
between the vertices in the two parts. There are $2^{k-1}-1$ such
partitions. Fix one, and fix a spanning tree $T$ of the complete bipartite
graph between its two parts. The tree $T$ has $k-1$ edges, and if no edge of
$H$ goes across the partition, every edge of $T$ joins two nonadjacent
vertices of $H$. Root $T$, and reveal the vertices of $S_i$ in an order of
the positions in which every position comes after its parent. When a
position other than the root is revealed, at least $n-k+1$ vertices are still
available, all equally likely, and at most $q$ of them are nonneighbors of
the vertex at the parent position. So the probability that every edge of $T$
joins nonadjacent vertices is at most $(q/(n-k+1))^{k-1}$. By the union bound
over the partitions, $H[S_i]$ is disconnected with probability at most $F$,
and the expected number of disconnected blocks is at most $NF$.

\emph{Missing pairs.} Let $i<j$. Condition on $S_i$, and suppose that $S_i$
is good. Given $S_i$, the block $S_j$ is a uniformly random $k$-subset of the
$n-k$ vertices outside $S_i$. If no edge joins $S_i$ and $S_j$, then
$S_j\subseteq U(S_i)$. Hence, as $|U(S_i)|\le n-k$,
\begin{equation*}
 \Prb\bigl(\text{no edge joins }S_i\text{ and }S_j\bigm|S_i\bigr)
 \le\binom{|U(S_i)|}{k}\Big/\binom{n-k}{k}
 \le\left(\frac{|U(S_i)|}{n-k}\right)^k
 \le\left(\frac{An\pi^k}{n-k}\right)^k.
\end{equation*}
Let $Z$ be the number of pairs $i<j$ such that $S_i$ and $S_j$ are both good
and connected and no edge joins them. Dropping the conditions that $S_j$ is
good and that both blocks are connected can only increase the probability of
such a pair. So the bound above, averaged over $S_i$ and summed over the
$\binom N2$ pairs, gives $\E Z\le M$.

\emph{The model.} Delete every block that is not good or is disconnected.
Then, for each pair of remaining blocks with no edge between them, delete one
of the two blocks. This deletes at most $Z$ further blocks. The remaining
blocks are disjoint, each induces a connected subgraph, and every two of them
are joined by an edge. So they form a $K_t$ model in $H$, where $t$ is their
number. By linearity of expectation, $\E t\ge N-N/A-NF-M=B$. Hence some
ordering gives $t\ge B>m-1$, and since $t$ is an integer, $t\ge m$. So $H$ has
a $K_m$ minor.
\end{proof}

The proof uses only the distribution of one block, or of one block given
another, and linearity of expectation. It does not need any independence
between the blocks.

\subsection{The best choice of \texorpdfstring{$A$}{A}}

Lemma~\ref{lem:partition} holds for every $A>0$. For $q>0$ we choose $A$ to
make $N/A+M$ as small as possible. Write $M=c_0A^k$, with
$c_0=\binom N2(n\pi^k/(n-k))^k$. As a function of $A>0$, the sum
$N/A+c_0A^k$ is smallest when $A^{k+1}=N/(kc_0)$, and then $c_0A^k=N/(kA)$, so
that $N/A+M=(1+1/k)N/A$. Put $P=(1+1/k)/A$ for this $A$. Then
Lemma~\ref{lem:partition} holds with
\begin{equation}\label{eq:optimized}
 B=N(1-F-P),\qquad
 P^{k+1}=\left(1+\frac1k\right)^{k+1}
 \frac{k(N-1)}2\left(\frac n{n-k}\right)^k\pi^{k^2}.
\end{equation}
This also holds for $q=0$, with $F=P=0$. Indeed, then $B=N(1-1/A)$ in
Lemma~\ref{lem:partition}, and if $N>m-1$, then $N(1-1/A)>m-1$ for every
large enough $A$.

\subsection{Blocks of sizes 2 and 3}

For some cases of Section~\ref{sec:finite} it helps to use blocks of two
sizes.

\begin{lemma}\label{lem:mixed}
Let $H$ be a graph with $n$ vertices in which every vertex has at most $q$
nonneighbors, and put $\pi=q/n$. Let $N_2,N_3\ge1$ be integers with
$2N_2+3N_3\le n$, and let $A_2,A_3>0$. Put
\begin{align*}
 F_2&=\frac q{n-1},& F_3&=3\left(\frac q{n-2}\right)^2,\\
 M_{22}&=\binom{N_2}{2}\left(\frac{A_2n}{n-2}\right)^2\pi^4,&
 M_{33}&=\binom{N_3}{2}\left(\frac{A_3n}{n-3}\right)^3\pi^9,\\
 M_{23}&=N_2N_3\left(\frac{A_2n}{n-2}\right)^3\pi^6,
\end{align*}
and
\begin{equation}\label{eq:mixed}
 B=N_2(1-F_2-A_2^{-1})+N_3(1-F_3-A_3^{-1})-M_{22}-M_{23}-M_{33}.
\end{equation}
If $B>m-1$ for an integer $m$, then $H$ has a $K_m$ minor.
\end{lemma}

\begin{proof}
We follow the proof of Lemma~\ref{lem:partition}. In a uniformly random
ordering of $V(H)$, the first $2N_2$ vertices form $N_2$ blocks of size $2$,
called pairs, and the next $3N_3$ vertices form $N_3$ blocks of size $3$,
called triples. A pair $S$ is good if $|U(S)|\le A_2n\pi^2$, and a triple
$S$ is good if $|U(S)|\le A_3n\pi^3$. The domination and connectivity bounds of
Lemma~\ref{lem:partition} with $k=2$ and $k=3$ show that the expected
numbers of pairs and triples that are not good or are disconnected are at
most $N_2(F_2+1/A_2)$ and $N_3(F_3+1/A_3)$. For two pairs, or for two
triples, the bound on missing pairs of Lemma~\ref{lem:partition} gives the
terms $M_{22}$ and $M_{33}$. For a pair and a triple, condition on the pair
and suppose that it is good. The triple is then a uniformly random
$3$-subset of the $n-2$ other vertices, and it has no neighbor in the pair
only if it lies in the set of vertices that the pair does not dominate. This
happens with probability at most $(A_2n\pi^2/(n-2))^3$, and summing over the
$N_2N_3$ choices of a pair and a triple gives $M_{23}$. The rest of the proof
is the same as for Lemma~\ref{lem:partition}.
\end{proof}

\section{Exact certificate for \texorpdfstring{$51\le p\le109$}{51 <= p <= 109}}\label{sec:finite}

In this range we check, for each $p$ and each order $n$ of a reduced graph,
that Lemma~\ref{lem:partition} or Lemma~\ref{lem:mixed} with $m=p-2$ applies.
Then every reduced graph has a $K_{p-2}$ minor.

The file \texttt{global\_finite\_certificate.jsonl} has one row for each
integer $p$ from $51$ to $109$. A row gives an integer $c$ and one choice of
blocks for each order
\begin{equation*}
 n=R+2,R+3,\ldots,c,\qquad\text{where }R=\floor{25c/67}-1.
\end{equation*}
The checker \texttt{global\_finite\_checker.py} first verifies, with a
rational lower bound $\ell_p$ on $\ln p$ (Section~\ref{sec:enclosures}), that
\begin{equation*}
 c^2<C^2p^2\ell_p\qquad\text{and}\qquad R\ge p-2.
\end{equation*}
The first inequality gives $c<D$. The second gives $c\ge3$, since it needs
$\floor{25c/67}\ge50$, and it handles
every complete reduced graph, since such a graph has at least $R\ge p-2$
vertices. By \eqref{eq:parameters}, this includes every reduced graph of
order $n=R$ or $n=R+1$. For each $n\ge R+2$ the checker uses the bound
\begin{equation*}
 q=\floor{\frac{n-R}{2}},
\end{equation*}
the largest number of nonneighbors of a vertex that a reduced graph of order
$n$ allows.
Lemmas~\ref{lem:partition} and~\ref{lem:mixed} need only an upper bound $q$
on the number of nonneighbors of each vertex, so a check for this value of
$q$ covers every reduced graph of order $n$. (Equivalently, every loss term
in \eqref{eq:B} and \eqref{eq:mixed} is nondecreasing in $q$.)

A uniform entry of a row gives an integer $k$ and a positive rational $A$,
for Lemma~\ref{lem:partition}. A mixed entry gives integers $N_2,N_3$ and
positive rationals $A_2,A_3$, for Lemma~\ref{lem:mixed}. For each entry the
checker verifies the conditions of the lemma, computes $B$ from
\eqref{eq:B} or \eqref{eq:mixed} with exact fractions, and verifies
$B>p-3$. It also verifies that the rows cover every $p$ from $51$ to $109$,
and that each row covers every $n$ from $R+2$ to $c$. The result is as
follows.
\begin{center}
\begin{tabular}{lr}
\toprule
Quantity & Verified value\\
\midrule
Values of $p$ & $51,52,\ldots,109$\\
Cases $(p,n)$ checked & $19\,873$\\
Cases using Lemma~\ref{lem:mixed} & $313$\\
Lower bound for every $B-(p-3)$ & $1/10$ (strict)\\
\bottomrule
\end{tabular}
\end{center}
The smallest margin $B-(p-3)$ is about $0.100642159415$. The checker verifies
with exact fractions that every margin is greater than $1/10$. So
Lemma~\ref{lem:partition} or Lemma~\ref{lem:mixed} gives a $K_{p-2}$ minor
in every reduced graph for $51\le p\le109$.

\section{Exact interval certificate for \texorpdfstring{$110\le p\le10^{44}$}{110 <= p <= 10\^{}44}}
\label{sec:medium}

For larger $p$ we cannot check each value of $p$ on its own. So the second
certificate covers whole intervals of $p$ at once. Within an interval, it
also covers the order $n$ of a reduced graph by a few ranges, after dividing
$n$ by $D$. Every check is again exact.

Fix integers $110\le L\le U$, and let $s_L$ and $s_U$ be rational numbers
with
\begin{equation*}
 s_L\le\sqrt{\ln L}\qquad\text{and}\qquad s_U\ge\sqrt{\ln U}.
\end{equation*}
Section~\ref{sec:enclosures} explains how they are computed. Let
$D_L=CLs_L$, $c_L=\ceil{D_L}-1$ and $R_L=\floor{25c_L/67}-1$. Now let $p$ be
an integer with $L\le p\le U$, and let $D$, $c$ and $R$ be as in \eqref{eq:D}
and \eqref{eq:R}. Then $CLs_L\le D\le CUs_U$, because
$s_L\le\sqrt{\ln L}\le\sqrt{\ln p}\le\sqrt{\ln U}\le s_U$, and $c\ge3$ since
$D>4$.

We first bound $R/D$ from below. Write $\bar\alpha=67/25$. Since
$\floor{y}>y-1$ for every real $y$, and $c\ge D-1$,
\begin{equation*}
 R>\frac{c}{\bar\alpha}-2\ge\frac D{\bar\alpha}-\left(\frac1{\bar\alpha}+2\right).
\end{equation*}
Dividing by $D\ge D_L$ gives the first bound below. Also $c\ge c_L$, because
$D\ge D_L$, and hence $R\ge R_L$. With $D\le CUs_U$, this gives the second
bound below when $R_L\ge0$. So
\begin{equation*}
 \frac RD\ge\rho:=\max\left\{
 \frac1{\bar\alpha}-\frac{1/\bar\alpha+2}{D_L},\
 \frac{R_L}{CUs_U}\right\}.
\end{equation*}
The checker verifies $\rho>0$. If the second term is the larger one, it is
positive, so $R_L>0$ and the second bound applies.

Let $H$ be a reduced graph with $n$ vertices, let $q$ be the largest number of
nonneighbors of a vertex of $H$, and put $\pi=q/n$ and $x=n/D$. By
\eqref{eq:parameters}, $R\le n\le c<D$, so $\rho\le x\le1$.

The certificate splits $[\rho,1]$ into intervals $[a,b]$ with rational ends.
We call an interval of $p$ together with one interval of $x$ a \emph{cell}.
For a cell put
\begin{equation*}
 n_-=CLs_La,\qquad n_+=CUs_Ub,\qquad
 \pi_+=\frac{1-\rho/b}{2}.
\end{equation*}
If $x\in[a,b]$, then $n_-\le n\le n_+$, since $n=xD$. Also $\pi\le\pi_+$,
because $q\le(n-R)/2$ by \eqref{eq:parameters} and
$R/n=(R/D)/x\ge\rho/b$. Each cell gives an integer $k\ge2$ and a rational
$P_+\ge0$. The checker verifies $n_->k$, $n_-/k\ge2$ and $0\le\pi_+<1/2$, so
that $2\le k<n$ and $N=\floor{n/k}\ge2$, as Lemma~\ref{lem:partition}
requires.

We use Lemma~\ref{lem:partition} in the form \eqref{eq:optimized}, which
also covers $q=0$. Since $q/(n-k+1)=\pi/(1-(k-1)/n)$, the connectivity loss
satisfies $F\le F_+$, where
\begin{equation*}
 F_+=(2^{k-1}-1)
 \left(\frac{\pi_+}{1-(k-1)/n_-}\right)^{k-1}.
\end{equation*}
In \eqref{eq:optimized} we have $k(N-1)\le n\le n_+$, and $n/(n-k)$ decreases
as $n$ increases, so $n/(n-k)\le n_-/(n_--k)$. Hence $P^{k+1}\le Y$, where
\begin{equation*}
 Y=\left(1+\frac1k\right)^{k+1}\frac{n_+}{2}
 \left(\frac{n_-}{n_--k}\right)^k\pi_+^{k^2}.
\end{equation*}
The checker verifies $P_+^{k+1}\ge Y$, which gives $P\le P_+$ without
computing any root. It also verifies $F_++P_+<1$ and
\begin{equation}\label{eq:celltest}
 \left(\frac{Cs_La}{k}-\frac1L\right)(1-F_+-P_+)
 >1-\frac3U.
\end{equation}

This gives $B>p-3$, as follows. Since $N>n/k-1$, $n\ge aD$ and
$D\ge Cps_L$, we have
\begin{equation*}
 \frac Np>\frac n{kp}-\frac1p\ge\frac{Cs_La}{k}-\frac1L.
\end{equation*}
As $1-F-P\ge1-F_+-P_+>0$, this and \eqref{eq:celltest} give
\begin{equation*}
 \frac Bp=\frac Np(1-F-P)
 >\left(\frac{Cs_La}{k}-\frac1L\right)(1-F_+-P_+)
 >1-\frac3U\ge1-\frac3p,
\end{equation*}
which is $B>p-3$.

The file \texttt{global\_medium\_certificate.jsonl} has $88$ intervals of $p$
and $508$ cells. The checker \texttt{global\_medium\_checker.py} verifies
that the first interval starts at $110$, that each interval starts one after
the end of the previous one, and that the last one ends at $10^{44}$. Within
each interval it verifies that the intervals of $x$ start at the computed
$\rho$, that each starts where the previous one ends, and that the last one
ends at $1$. Every cell passes \eqref{eq:celltest}. The smallest margin of
the left side of \eqref{eq:celltest} above $1-3/U$ is about
$9.98234667206\times10^{-7}$, and the checker verifies with exact fractions
that every margin is positive.

\subsection{Exact bounds on logarithms and square roots}\label{sec:enclosures}

For a rational number $z$ with $0\le z<1$ and for $J=24$, put
\begin{equation*}
 T_J(z)=2\sum_{j=0}^{J-1}\frac{z^{2j+1}}{2j+1}.
\end{equation*}
The series $2\atanh z=\sum_{j\ge0}2z^{2j+1}/(2j+1)$ has positive terms, and
its tail after $J$ terms is at most the geometric series
$\sum_{j\ge J}2z^{2j+1}/(2J+1)$. So
\begin{equation}\label{eq:atanh}
 T_J(z)\le2\atanh z
 \le T_J(z)+\frac{2z^{2J+1}}{(2J+1)(1-z^2)}.
\end{equation}
For an integer $m\ge1$ let $h=\floor{\log_2m}$, which is one less than the
number of binary digits of $m$, and put $z=(m-2^h)/(m+2^h)$. Then
$0\le z<1/3$, and
\begin{equation*}
 \ln m=h\ln2+2\atanh z,
 \qquad \ln2=2\atanh(1/3),
\end{equation*}
because $2\atanh z=\ln((1+z)/(1-z))$. Applying \eqref{eq:atanh} to both
terms gives rational bounds $\ell_m\le\ln m\le u_m$. The finite checker uses
only the lower bound.

For the interval checker, let $S=10^{12}$, and compute with integer square
roots
\begin{equation*}
 s_m^- =\frac{\floor{\sqrt{\floor{S^2\ell_m}}}}S,
 \qquad
 s_m^+ =\frac{\floor{\sqrt{\floor{S^2u_m}}}+1}S.
\end{equation*}
Then $s_m^-\le\sqrt{\ell_m}\le\sqrt{\ln m}$, and $s_m^+>\sqrt{u_m}\ge\sqrt{\ln
m}$ because $(Ss_m^+)^2\ge\floor{S^2u_m}+1>S^2u_m$. The checker takes
$s_L=s_L^-$ and $s_U=s_U^+$. All quantities that the checkers accept are
computed with integer or rational arithmetic. The scripts also contain
generators, which use floating point to propose certificates. The checkers do
not rely on the generators.

\section{An analytic proof for \texorpdfstring{$\ln p\ge100$}{ln p >= 100}}\label{sec:tail}

For very large $p$ we show $B>p-3$ by hand, with $A=4$. The block size $k$
grows like $\sqrt{\ln p}$, which makes the connectivity loss and the
missing-pair loss small. Put $s=\sqrt{\ln p}\ge10$, so that $p=e^{s^2}$. Let
$H$ be a reduced graph with $n$ vertices, and let $q$ be the largest number of
nonneighbors of a vertex of $H$. With $c$ and $R$ from \eqref{eq:D} and
\eqref{eq:R},
\begin{equation}\label{eq:tailR}
 R\ge\frac{80}{67}ps-\frac{159}{67}\ge1.19ps,
 \qquad n\le3.2ps.
\end{equation}
The first inequality holds because $R>25c/67-2$ and $c\ge D-1$. The second
follows from $80/67-119/100=27/6700$ and $ps\ge600$, which holds as
$p\ge e^{100}$. The bound on $n$ holds
because $n\le c<D$.

If $q=0$, then $H$ is complete and $n\ge R>p-2$, so $H$ has a $K_{p-2}$
minor. Suppose $q>0$, and let
\begin{equation*}
 \pi=q/n\qquad\text{and}\qquad\lambda=\ln(1/\pi).
\end{equation*}
By \eqref{eq:parameters} and \eqref{eq:tailR},
\begin{equation*}
 0<\pi\le\frac{1-R/n}{2}\le\frac{1-1.19/3.2}{2}=\frac{201}{640}<\frac13,
 \qquad\text{so}\qquad \lambda>\ln3>1.09.
\end{equation*}
Here $\ln3=2\atanh(1/2)>1+1/12+1/80+1/448=7379/6720>1.09$, from the first
four terms of the series. Choose
\begin{equation*}
 k=\ceil{\frac{s}{\sqrt\lambda}+1},\qquad
 N=\floor{n/k},\qquad A=4.
\end{equation*}
Since $\lambda>1$, we have $2\le k<s/\sqrt\lambda+2<s+2$. So, by
\eqref{eq:tailR} and $s/(s+2)\ge10/12$,
\begin{equation*}
 \frac nk>\frac{1.19ps}{s+2}\ge\frac{119}{120}p>1000.
\end{equation*}
In particular $N\ge2$, $k/n<1/1000$ and $n/(n-k)<1000/999$. As
$e^{7/10}>\sum_{j=0}^3(7/10)^j/j!>2000/999$ and
$e^{7/5}>\sum_{j=0}^5(7/5)^j/j!>4000/999$, we get
\begin{equation}\label{eq:logs}
 \ln\frac{2n}{n-k+1}<0.7,
 \qquad \ln\frac{4n}{n-k}<1.4.
\end{equation}

\subsection{Connectivity loss}

The loss $F$ of Lemma~\ref{lem:partition} satisfies
\begin{align*}
 F&<\left(\frac{2q}{n-k+1}\right)^{k-1}
 <\exp(-(k-1)(\lambda-0.7))\\
 &\le\exp\left[-s\left(\sqrt\lambda-\frac{0.7}{\sqrt\lambda}\right)\right]
 < e^{-0.37s}\le e^{-3.7}<0.03.
\end{align*}
The first step uses $2^{k-1}-1<2^{k-1}$. The second uses
$2q/(n-k+1)=\pi\cdot2n/(n-k+1)$ and \eqref{eq:logs}. The third uses
$k-1\ge s/\sqrt\lambda$ and $\lambda>0.7$. For the fourth, the function
$\sqrt\lambda-0.7/\sqrt\lambda$ increases with $\lambda$, and at
$\lambda=1.09$ it equals $0.39/\sqrt{1.09}>0.37$, since
$0.39^2=0.1521>0.149221=0.37^2\cdot1.09$. The last two steps use $s\ge10$ and
$e^{3.7}>\sum_{j=0}^5 3.7^j/j!>100/3$.

\subsection{Missing-pair loss}

With $A=4$, the term $M$ of Lemma~\ref{lem:partition} and \eqref{eq:logs}
give
\begin{equation*}
 \frac MN=\frac{N-1}{2}\left(\frac{4n\pi^k}{n-k}\right)^k
 \le\frac{N-1}{2}\exp(1.4k-k^2\lambda)
 \le1.6ps\exp(1.4k-k^2\lambda),
\end{equation*}
where the last step uses $N-1<n\le3.2ps$. The function $y^2\lambda-1.4y$ is
increasing for $y\ge0.7/\lambda$, and $k\ge s/\sqrt\lambda+1>0.7/\lambda$.
Hence
\begin{align*}
 k^2\lambda-1.4k
 &\ge\left(\frac{s}{\sqrt\lambda}+1\right)^2\lambda
   -1.4\left(\frac{s}{\sqrt\lambda}+1\right)
 =s^2+\left(2\sqrt\lambda-\frac{1.4}{\sqrt\lambda}\right)s
       +\lambda-1.4\\
 &>s^2+0.74s-0.31,
\end{align*}
where the last step uses
$2\sqrt\lambda-1.4/\sqrt\lambda>2\cdot0.37$ from the previous subsection and
$\lambda>1.09$. Since $p=e^{s^2}$,
\begin{equation*}
 \frac MN<1.6s\exp(-0.74s+0.31)<0.02.
\end{equation*}
The middle expression decreases for $s\ge10$, since the derivative of
$se^{-0.74s}$ is $(1-0.74s)e^{-0.74s}$. At $s=10$ it equals
$16e^{-7.09}$, which is less than $0.02$ because
$e^{7.09}>\sum_{j=0}^87.09^j/j!>800$. Combining the two losses with
$A^{-1}=1/4$ gives
\begin{equation}\label{eq:surviving}
 B>N(1-0.03-0.25-0.02)=0.70N.
\end{equation}

\subsection{There are enough blocks}

For $0<\pi\le1/2$,
\begin{equation*}
 \frac{\dd}{\dd\pi}\left[\ln(1/\pi)-2(1-2\pi)^2\right]
 =-\frac1\pi+8(1-2\pi)=-\frac{(4\pi-1)^2}{\pi}\le0.
\end{equation*}
So the bracket is at least its value $\ln2$ at $\pi=1/2$, and
\begin{equation*}
 \lambda\ge2(1-2\pi)^2+\ln2
 \ge(2+\ln2)(1-2\pi)^2,
\end{equation*}
where the last step uses $(1-2\pi)^2\le1$. The first two terms of the series
for $2\atanh(1/3)$ give $\ln2>56/81>2/3$, and $0.613^2=0.375769>3/8$. Hence
\begin{equation*}
 \frac{1-2\pi}{\sqrt\lambda}
 \le\frac1{\sqrt{2+\ln2}}<\sqrt{3/8}<0.613.
\end{equation*}
Now $n(1-2\pi)=n-2q\ge R\ge1.19ps$ by \eqref{eq:parameters} and
\eqref{eq:tailR}, and $k<s/\sqrt\lambda+2$. So
\begin{equation*}
 \frac nk
 >\frac{1.19ps}{(1-2\pi)(s/\sqrt\lambda+2)}
 =\frac{1.19p}{(1-2\pi)/\sqrt\lambda+2(1-2\pi)/s}
 \ge\frac{1.19p}{0.613+2/s}
 \ge\frac{1.19p}{0.813}>1.46p.
\end{equation*}
Thus $N>n/k-1>1.46p-1$, and \eqref{eq:surviving} gives
\begin{equation*}
 B>0.70(1.46p-1)=1.022p-0.7>p-3.
\end{equation*}
So Lemma~\ref{lem:partition} with $m=p-2$ gives a $K_{p-2}$ minor in $H$.

\section{Proof of the Main Theorem}\label{sec:proof}

Let $p\ge51$. By Section~\ref{sec:reduction}, it suffices to show that every
reduced graph has a $K_{p-2}$ minor, for the $c$ of \eqref{eq:D}, or for the
$c$ of the certificate when $p\le109$. The three ranges are covered as
follows.
\begin{center}
\begin{tabular}{ll}
\toprule
Range & Argument\\
\midrule
$51\le p\le109$ & $19\,873$ exact checks (Section~\ref{sec:finite})\\
$110\le p\le10^{44}$ & $88$ intervals of $p$ and $508$ exact cells (Section~\ref{sec:medium})\\
$\ln p\ge100$ & analytic argument (Section~\ref{sec:tail})\\
\bottomrule
\end{tabular}
\end{center}
The interval checker also verifies $\ell_m>100$ for $m=10^{44}$, with the
lower bound of Section~\ref{sec:enclosures}. So $\ln p>100$ for every
$p>10^{44}$, and the three ranges cover every $p\ge51$. In each case
Lemma~\ref{lem:partition} or Lemma~\ref{lem:mixed} gives a $K_{p-2}$ model in
every reduced graph, and Thomason's lemma gives the Main Theorem. Together
with the Lean proof for $p\le50$ (Section~\ref{sec:lean}), this proves (5.5)
of Graph Minors XIII for every $p\ge1$, with no unspecified threshold on $p$.

The two checkers and their certificates are in the folder
\texttt{docs/thomason} of the GraphMinor repository
(\url{https://github.com/mahdi-jfri/GraphMinor}), next to this note. They
need Python~3.8 or later and no other package. The commands
\begin{verbatim}
python3 global_finite_checker.py check global_finite_certificate.jsonl
python3 global_medium_checker.py check global_medium_certificate.jsonl
\end{verbatim}
run the checks of Sections~\ref{sec:finite} and~\ref{sec:medium}, and
\texttt{python3 verify\_all.py} runs both. The certificate files contain only
witnesses, that is, the choices of blocks and parameters. The checkers verify
the coverage of the ranges and every inequality that the proof uses.

\section{The Lean formalization}\label{sec:lean}

The Lean formalization proves (5.5) of Graph Minors XIII for every $p\ge1$,
as the theorem \texttt{GraphMinor.Thomason.thomason84} in
\texttt{GraphMinor/Thomason/UniformTheorem.lean}. It splits the proof into
four ranges of $p$. All file names below are in the folder
\texttt{GraphMinor/Thomason}.

\emph{The range $p\le50$.} The proof is in \texttt{FiftyCase.lean} and the
files it imports. It does not follow this note.

\emph{The density reduction.} Thomason's lemma is proved in
\texttt{DensityReduction.lean}, following \cite{thomason}, for every real
$a>1$ with $1+\ln(2a)\le a$ in place of $\alpha_0$. The file
\texttt{UniformReduction.lean} applies it with $a=67/25$, which satisfies
this condition by the bound $e^{42/25}>134/25$ of
Section~\ref{sec:reduction}. It gives the reduction to reduced graphs of
Section~\ref{sec:reduction} with $R=\floor{25c/67}-1$.

\emph{The random partition.} Lemma~\ref{lem:partition} is proved in
\texttt{UniformPartition.lean}, in a form that allows blocks of different
sizes. The expectations of the proof are written as averages over all
injective maps from the positions of the blocks to the vertices
(\texttt{UniformSamples.lean}, \texttt{UniformBlockCounts.lean},
\texttt{UniformBlockSums.lean} and \texttt{UniformPairSums.lean}). The file
\texttt{UniformTwoSizes.lean} specializes it to blocks of two sizes, which
covers Lemma~\ref{lem:partition} and Lemma~\ref{lem:mixed}, and
\texttt{UniformCheck.lean} turns its hypothesis into a test on natural
numbers that Lean's kernel can evaluate.

\emph{The range $51\le p\le109$.} The Lean proof uses the certificate of
Section~\ref{sec:finite}. The script \texttt{gen\_lean\_finite.py} copies the
value of $c$ and the block choices of each row of
\texttt{global\_finite\_certificate.jsonl} unchanged into the files
\texttt{UniformFinite51.lean} to \texttt{UniformFinite109.lean}, and the
kernel checks each row (\texttt{rowOK} in \texttt{UniformRows.lean}). The
only difference is the density test: the kernel checks
$c\le2.68\,p\sqrt{\log_2p}$ directly, with integers $a$ and $b=1000$ such that
$2^a\le p^b$ and $10000\,bc^2\le71824\,ap^2$, instead of $c<D$. The file
\texttt{UniformFiniteRange.lean} collects the rows.

\emph{The range $110\le p\le10^{44}$.} Here the Lean proof uses its own
certificate in base $2$, not \texttt{global\_medium\_certificate.jsonl}. It
keeps the density $D=2.68\,p\sqrt{\log_2p}$ of (5.5) itself, with
$c=\floor{D}$, and follows the method of Section~\ref{sec:medium}. The
logarithms are bounded by the integer comparisons $2^{a_L}\le L^{1000}$ and
$U^{1000}\le2^{a_U}$, and each cell also gives a rational upper bound on
$\pi$. Instead of the best $A$ of \eqref{eq:optimized}, each cell uses
$A=(k+1)/(kP_+)$, which gives the same bound on the losses and needs no
root. The certificate is \texttt{medium\_log2\_certificate.jsonl}, with $69$
intervals and $261$ cells, made and checked by
\texttt{medium\_log2\_certificate.py}. The script
\texttt{gen\_lean\_medium.py} writes it into the files
\texttt{UniformMedium0.lean} to \texttt{UniformMedium6.lean} and
\texttt{UniformMediumRange.lean}. The kernel checks each interval
(\texttt{intervalOK} in \texttt{UniformIntervalCheck.lean}), and
\texttt{UniformCell.lean} and \texttt{UniformIntervals.lean} prove that an
accepted interval gives a $K_{p-2}$ minor in every reduced graph.

\emph{The range $p>10^{44}$.} The files \texttt{UniformTail.lean} and
\texttt{UniformLarge.lean} replace Section~\ref{sec:tail} by an argument in
base $2$. With $s=\sqrt{\log_2p}\ge12$, so that $p=2^{s^2}$, they split the
range of $\pi=q/n$ into the three pieces $\pi\le1/25$,
$1/25<\pi\le13/100$ and $13/100<\pi\le63/200$. On each piece they take
blocks of size $\ceil{ts}+1$, for a fixed $t$ that depends on the piece, and
$A=4$. The connectivity loss and the missing-pair loss per block are then
each at most $1/50$, and there are more than $1.41p-1$ blocks, so
Lemma~\ref{lem:partition} applies with
$B>0.71(1.41p-1)>p-3$.

\appendix
\section{Proof of Thomason's lemma}\label{app:reduction}

Let $p\ge2$ and $c\ge3$ be integers, put $r=c/\alpha_0$, and suppose that
every non-null graph satisfying \eqref{eq:thomason} has a $K_{p-2}$ minor.
Since $\alpha_0<67/25$, we have $r>1$. Let $G$ be a non-null graph with
$e(G)\ge c\,v(G)$. We find a $K_p$ minor in two steps. Each step passes to a
minimal minor and then to the neighborhood of a vertex of minimum degree.
That vertex then gives one more branch set.

\emph{Step 1.} Among the non-null minors $G_1$ of $G$ with
$e(G_1)\ge c\,v(G_1)$, choose one with $v(G_1)+e(G_1)$ as small as possible.
Then no deletion of a vertex, deletion of an edge, or contraction of an edge of
$G_1$ gives another such minor. We have $e(G_1)=c\,v(G_1)$, since otherwise
$e(G_1)\ge c\,v(G_1)+1$, as $c$ is an integer, and deleting an edge would keep
the inequality. As $c\,v(G_1)=e(G_1)\le v(G_1)(v(G_1)-1)/2$, we have
$v(G_1)\ge2c+1$. Let $z$ be a vertex of minimum degree in $G_1$, and let
$H_2=G_1[N(z)]$. Then $v(H_2)=\deg z\le2e(G_1)/v(G_1)=2c$. Deleting $z$ leaves
a graph with $v(G_1)-1$ vertices and $e(G_1)-\deg z$ edges, so minimality
gives $\deg z\ge c+1$. For each neighbor $w$ of $z$, contracting $zw$ removes
one vertex and exactly $1+|N(z)\cap N(w)|$ edges, so minimality gives
$e(G_1)-1-|N(z)\cap N(w)|<c(v(G_1)-1)$, that is, $|N(z)\cap N(w)|\ge c$. So
every vertex of $H_2$ has degree at least $c$ in $H_2$, and
\begin{equation*}
 c+1\le v(H_2)\le2c,\qquad \delta(H_2)\ge c.
\end{equation*}
Since $z$ is adjacent to every vertex of $H_2$, a $K_{p-1}$ model in $H_2$
together with $\{z\}$ is a $K_p$ model in $G_1$.

\emph{Step 2.} For real $x\ge r$ put
\begin{equation*}
 \Psi(x)=\frac{rx}{2}\left(\ln\frac xr+1\right).
\end{equation*}
This function is increasing, since its derivative is
$(r/2)(\ln(x/r)+2)>0$. For $r\le x\le2c$ we have $\Psi(x)\le cx/2$, because
$\ln(x/r)\le\ln(2\alpha_0)=\alpha_0-1$ by the definition of $\alpha_0$, so
that $r(\ln(x/r)+1)\le\alpha_0r=c$. As $r<c+1\le v(H_2)\le2c$, this gives
$\Psi(v(H_2))\le c\,v(H_2)/2\le e(H_2)$, where the last step uses
$\delta(H_2)\ge c$.

Among the minors $J$ of $H_2$ with $v(J)\ge r$ and $e(J)\ge\Psi(v(J))$,
choose one with $v(J)+e(J)$ as small as possible, and put $n=v(J)$. Such
minors exist, since $H_2$ itself is one. No graph
with $r\le n<r+1$ has these properties, because then
\begin{equation*}
 \binom n2<\frac{rn}{2}\le\Psi(n).
\end{equation*}
So $n\ge r+1$, and deleting a vertex or contracting an edge of $J$ leaves a
graph with at least $r$ vertices. Since $e(J)$ is an integer and deleting an
edge does not change $n$, minimality gives $e(J)=\ceil{\Psi(n)}$. The graph $J$ has no isolated vertex, because
deleting one would keep $e(J)\ge\Psi(n)\ge\Psi(n-1)$, as $\Psi$ is
increasing.

Let $u$ be a vertex of minimum degree in $J$, and let $H=J[N_J(u)]$. Then $H$
is non-null, since $u$ is not isolated. Since $n\ge r+1>2$ and
$n\le v(H_2)\le2c$,
\begin{equation*}
 v(H)=\deg u\le\frac{2\ceil{\Psi(n)}}n
 <\frac{2\Psi(n)+2}{n}\le r\left(\ln\frac nr+1\right)+1.
\end{equation*}
As $v(H)$ is an integer, this gives
\begin{equation*}
 v(H)\le\ceil{r\left(\ln\frac nr+1\right)}
 \le\ceil{\alpha_0r}=c,
\end{equation*}
where the second step uses $n\le2c$, as in the bound $\Psi(x)\le cx/2$ above.
For each neighbor $w$ of $u$, contracting $uw$ removes one vertex and
$1+|N_J(u)\cap N_J(w)|$ edges, so minimality gives
\begin{equation*}
 e(J)-1-|N_J(u)\cap N_J(w)|<\Psi(n-1).
\end{equation*}
The degree of $w$ in $H$ is $|N_J(u)\cap N_J(w)|$, so
$\delta(H)>\Psi(n)-\Psi(n-1)-1$. With the bound on $v(H)$ above,
\begin{align*}
 2\delta(H)-v(H)
 &>2\Psi(n)-2\Psi(n-1)-r\left(\ln\frac nr+1\right)-3\\
 &=r(n-1)\ln\frac n{n-1}-3>r-4.
\end{align*}
The equality follows from the definition of $\Psi$. The last step uses
$\ln(n/(n-1))>1/n$, which gives $r(n-1)\ln(n/(n-1))>r-r/n>r-1$, as $n>r$.
Since $2\delta(H)-v(H)$ is an integer, it is at least
$\floor r-3=\floor{c/\alpha_0}-3$. So $H$ satisfies \eqref{eq:thomason}, and
by assumption $H$ has a $K_{p-2}$ model. Adding $\{u\}$ to it gives a
$K_{p-1}$ model in $J$. Since $J$ is a minor of $H_2$, the graph $H_2$ has a
$K_{p-1}$ minor, and by Step~1 the graph $G_1$ has a $K_p$ minor. As $G_1$ is
a minor of $G$, so does $G$.

\begin{thebibliography}{RS95}
\bibitem[RS95]{gm13}
N.~Robertson and P.~D.~Seymour,
\emph{Graph Minors. XIII. The disjoint paths problem},
J. Combin. Theory Ser. B \textbf{63} (1995), 65--110.
\href{https://doi.org/10.1006/jctb.1995.1006}{doi:10.1006/jctb.1995.1006}.
The statement (5.5) is on page~83.

\bibitem[Tho84]{thomason}
A.~Thomason, \emph{An extremal function for contractions of graphs},
Math. Proc. Cambridge Philos. Soc. \textbf{95} (1984), 261--265.
\href{https://doi.org/10.1017/S0305004100061521}{doi:10.1017/S0305004100061521}.
The density reduction is Lemma~1 on page~262, and the asymptotic conclusion
is on page~264.

\bibitem[Woo26]{wood}
D.~R.~Wood, \emph{An explicit bound on the extremal function for complete
graph minors}, arXiv:2610.41060v1, 29 October 2026.
\url{https://arxiv.org/abs/2610.41060v1}.
The bound is Theorem~1.
\end{thebibliography}
\end{document}
